% ============================================================
% paper/appendices/appx_repro_manifest_pointer.tex
% Appendix A: Reproducibility Manifest Pointer
% ============================================================

\section*{Appendix A: Reproducibility Manifest Pointer}

\subsection*{Purpose of This Appendix}

This appendix provides a formal pointer to reproducibility artifacts
associated with the Enterprise Continuum $K_{\mathrm{EA}}$.
It introduces no additional theory, methodology, assumptions, or
executable logic.

Its sole purpose is to delimit the relationship between this
whitepaper, the Executable Theory Specification, and any external
artifacts that may be used for structural reconstruction, inspection,
or audit.

\subsection*{Executable Source of Truth}

The definitive source of truth for the Enterprise Continuum is the
immutable Executable Theory Specification:

\begin{center}
\texttt{ETS.K\_EA.v1.1}
\end{center}

as referenced in the front matter of this document.
All admissible primitives, relations, thresholds, invariants, phase
precedence rules, and termination conditions are fully specified
therein.

This whitepaper is a human-readable projection of that specification
and does not modify, extend, or reinterpret it.

\subsection*{Reproducibility Scope}

Reproducibility in the context of $K_{\mathrm{EA}}$ is structural rather
than empirical.
It refers exclusively to the ability to reconstruct the same formal
object from the same executable specification, without interpretation,
extension, or invention.

External artifacts such as configuration files, manifests, scripts, or
instantiation inputs may exist to support reconstruction, inspection,
or audit.
Such artifacts are not part of the theory and carry no independent
normative, diagnostic, or explanatory authority.

\subsection*{Non-Claims and Boundaries}

This appendix explicitly makes no claims regarding:
\begin{itemize}
  \item correctness, completeness, or reliability of external tooling,
  \item empirical validity of any instantiation or simulation,
  \item performance, accuracy, or predictive value of derived outputs,
  \item suitability of artifacts for operational, managerial, or
        decision-making use.
\end{itemize}

The presence, absence, or quality of external reproducibility artifacts
does not alter the admissibility conditions, invariants, No-Go results,
or STOP logic defined by the Executable Theory Specification.

\subsection*{Interpretation Rule}

If any external artifact, reconstruction attempt, simulation output, or
visual material appears to contradict the formal results implied by
\texttt{ETS.K\_EA.v1.1}, the executable specification takes precedence.

Beyond this formal pointer, the appendix deliberately remains silent.

%%%End of Document%%%
